\documentclass[12pt, a4paper]{article}

% === Pacchetti ===
\usepackage[utf8]{inputenc}
\usepackage[T1]{fontenc}
\usepackage[italian]{babel}
\usepackage{geometry}
\usepackage{booktabs}
\usepackage{array}
\usepackage{xcolor}
\usepackage{colortbl}
\usepackage{enumitem}
\usepackage{fancyhdr}
\usepackage{titlesec}
\usepackage{microtype}
\usepackage{hyperref}
\usepackage{mdframed}
\usepackage{tcolorbox}
\usepackage{parskip}
\usepackage{lmodern}
\usepackage{amsmath}


% === Geometria pagina ===
\geometry{
  top=2.5cm, bottom=2.5cm,
  left=3cm,  right=2.5cm,
  headheight=15pt
}

% === Palette colori ===
\definecolor{primaryblue}{HTML}{1A3A5C}
\definecolor{accentblue}{HTML}{2E7EBA}
\definecolor{lightgray}{HTML}{F5F7FA}
\definecolor{medgray}{HTML}{AABCD0}
\definecolor{passgreen}{HTML}{2E7D32}
\definecolor{failred}{HTML}{C62828}
\definecolor{warnAmber}{HTML}{E65100}
\definecolor{roweven}{HTML}{EFF4F9}
\definecolor{afternoonblue}{HTML}{EAF2FB}

% === Hyperref ===
\hypersetup{
  colorlinks=true,
  linkcolor=accentblue,
  urlcolor=accentblue,
  pdftitle={Report Sessione 07 Giugno 2026 (Pomeriggio) -- GraphRAG MTB},
  pdfauthor={Paolo},
  pdfsubject={Tesi Magistrale},
}

% === Header / Footer ===
\pagestyle{fancy}
\fancyhf{}
\fancyhead[L]{\small\color{medgray}\textit{Tesi Magistrale: GraphRAG per Molecular Tumor Board}}
\fancyhead[R]{\small\color{medgray}07 Giugno 2026}
\fancyfoot[C]{\small\color{medgray}\thepage}
\renewcommand{\headrulewidth}{0.4pt}
\renewcommand{\footrulewidth}{0pt}

% === Titoli sezione ===
\titleformat{\section}
{\Large\bfseries\color{primaryblue}}
{\thesection.}{0.6em}{}
[\titlerule]
\titleformat{\subsection}
{\normalsize\bfseries\color{accentblue}}
{\thesubsection.}{0.5em}{}
\titlespacing*{\section}{0pt}{1.4em}{0.6em}
\titlespacing*{\subsection}{0pt}{1em}{0.3em}

% === Box colorati ===
\tcbuselibrary{skins,breakable}
\newtcolorbox{summarybox}{
  enhanced, breakable,
  colback=lightgray, colframe=primaryblue,
  boxrule=1.2pt, arc=4pt,
  fonttitle=\bfseries\color{white},
  title={Stato Finale dell'Implementazione},
}
\newtcolorbox{insightbox}{
  enhanced, breakable,
  colback=afternoonblue, colframe=accentblue,
  boxrule=1pt, arc=4pt,
  fonttitle=\bfseries\color{white},
  title={Insight Metodologici ed Evidenze},
}
\newtcolorbox{keyresultbox}{
  enhanced, breakable,
  colback=lightgray, colframe=passgreen,
  boxrule=1.2pt, arc=4pt,
  fonttitle=\bfseries\color{white},
  title={Risultato Chiave},
}

% === Macro colori ===
\newcommand{\pass}[1]{{\color{passgreen}\textbf{#1}}}
\newcommand{\fail}[1]{{\color{failred}\textbf{#1}}}
\newcommand{\warn}[1]{{\color{warnAmber}\textbf{#1}}}
\newcommand{\bench}[1]{{\texttt{\textbf{#1}}}}

% ============================================================
\begin{document}

% === Titolo ===
\begin{center}
  {\LARGE\bfseries\color{primaryblue} Report Sessione di Lavoro --- 07 Giugno 2026}\\[0.4em]
  {\large\color{accentblue} Progetto: Pipeline Agentica GraphRAG per Molecular Tumor Board (Tesi Magistrale)}\\[0.2em]
  {\small\color{medgray}\textit{Sessione pomeridiana (Definitivo)}}
\end{center}
\vspace{0.5em}\hrule\vspace{1.2em}

% ============================================================
\section{Punto di Partenza}
La sessione costituisce il proseguimento diretto delle attività concluse il 06/06/2026. All'avvio, il progetto presentava i seguenti elementi di partenza:
\begin{itemize}[leftmargin=1.6em, itemsep=3pt]
  \item \textbf{Sistema completo end-to-end} attivo con backend in FastAPI, orchestratore a nodi LangGraph, e frontend in React.
  \item \textbf{Benchmark iniziale su 30 casi} clinici con uno score medio complessivo di \textbf{3.96 / 5.0} calcolato su 5 criteri (incluso il parametro struttura).
  \item \textbf{Analisi critica del modulo judge:} identificazione di un'incoerenza nel criterio ``struttura'', che produceva sistematicamente un punteggio fisso di 5.0 su tutti i casi, gonfiando in modo fittizio la media complessiva.
  \item File \texttt{benchmark\_results.json} (versione v1) disponibile per l'analisi comparativa dei dati.
\end{itemize}

\textbf{Obiettivi principali della sessione:}
\begin{enumerate}[leftmargin=1.6em, itemsep=2pt]
  \item Ricalibrare e riorganizzare i criteri di valutazione del modulo \textit{LLM-as-judge} per allinearli alla letteratura scientifica.
  \item Ricalcolare i punteggi dell'intero benchmark di 30 casi clinici con la nuova griglia di valutazione.
  \item Analizzare in profondità i risultati e le metriche risultanti.
  \item Valutare l'efficacia clinica dell'OncoKB Enricher sui casi del benchmark con gap di copertura locale.
\end{enumerate}

% ============================================================
\section{Ricalibrazione Criteri LLM-as-Judge}

\subsection{Motivazione}
Il criterio ``struttura'' precedentemente utilizzato manteneva un valore costante di 5.0 in tutti i 30 casi del benchmark poiché il modulo \textit{Synthesizer} opera secondo istruzioni di prompt estremamente rigide che impongono una formattazione fissa. Questo indicatore non forniva alcuna misurazione reale dell'utilità clinica o dell'accuratezza scientifica e falsava la valutazione del sistema. Inoltre, i 5 criteri originari necessitavano di una formulazione più solida per risultare accademici e difendibili in sede di discussione di tesi.

\subsection{Nuovo Sistema a 4 Criteri (Allineati alla Letteratura)}
I parametri di valutazione sono stati riformulati ispirandosi direttamente a metodologie presenti in letteratura:
\begin{enumerate}[leftmargin=1.6em, itemsep=4pt]
  \item \textbf{Completezza (Comprehensiveness)} \textit{[Ispirato a Edge et al., 2025]}: Il report descrive esaustivamente tutti gli aspetti fondamentali del caso (variante genica, terapia indicata, profili di resistenza e trial clinici disponibili)?
  \item \textbf{Utilità Clinica (Clinical Actionability)} \textit{[Ispirato a Edge et al., 2025]}: Valuta l'immediata usabilità del report da parte degli specialisti del board. La correttezza formale e la leggibilità strutturale sono assorbite in questa dimensione, partendo dal presupposto che un report scarsamente organizzato compromette la sua stessa actionability.
  \item \textbf{Fedeltà alle Evidenze (Evidence Faithfulness)} \textit{[Ispirato a Wu et al., 2024]}: Verifica che i codici PubMed (PMID) citati siano realmente presenti nel Knowledge Graph locale e pertinenti alla raccomandazione formulata, eliminando fenomeni di allucinazione bibliografica.
  \item \textbf{Accuratezza Clinica (Clinical Accuracy)} \textit{[Criterio originale]}: Valuta la precisione scientifica delle indicazioni farmacologiche, delle companion diagnostics (CDx) associate e della corretta attribuzione del livello di evidenza ESCAT.
\end{enumerate}
Lo \textbf{Score Totale} viene ora calcolato come media aritmetica semplice dei 4 parametri sopra citati (ciascuno espresso in scala da 1.0 a 5.0).

\begin{insightbox}
\textbf{Citazione per la Tesi:}
\\
\textit{``I criteri di valutazione del Molecular Tumor Board Assistant sono stati ridefiniti allineando la metodologia LLM-as-judge alla letteratura scientifica recente, in particolare attingendo al framework di Edge et al. (2025) per la completezza e l'actionability, e all'approccio di Wu et al. (2024) per la precisione citazionale (Evidence Faithfulness), integrati con un criterio specifico di accuratezza clinica ad hoc per il tiering ESCAT.''}
\end{insightbox}

\subsection{Modifiche Codice}
Il nuovo assetto è stato implementato modificando:
\begin{itemize}[leftmargin=1.6em, itemsep=2pt]
  \item \texttt{backend/pipeline/agents/judge.py}: aggiornato il prompt \texttt{JUDGE\_SYSTEM} e lo schema JSON atteso.
  \item \texttt{backend/api/schemas.py}: aggiornato il modello Pydantic \texttt{JudgeResponse}.
  \item \texttt{frontend/src/types.ts}: modificata la corrispondente interfaccia TypeScript.
  \item \texttt{frontend/src/components/JudgePanel.tsx}: aggiornato il rendering grafico con le nuove etichette bilingue e i riferimenti accademici.
\end{itemize}

% ============================================================
\section{Benchmark Ricalcolato --- 30 Casi (v2)}

\subsection{Risultati Principali}
L'esecuzione sistematica della valutazione del benchmark con la nuova griglia a 4 criteri ha prodotto i seguenti risultati:
\begin{itemize}[leftmargin=1.6em, itemsep=3pt]
  \item \textbf{Casi totali valutati:} 30
  \item \textbf{Score medio generale:} \textbf{3.96 / 5.0} (stabile rispetto alla versione v1).
  \item \textbf{Score massimo registrato:} \textbf{5.00} (\bench{BENCH-029} BRAF V600E in carcinoma tiroideo).
  \item \textbf{Score minimo registrato:} \textbf{1.75} (\bench{BENCH-013} ALK G1202R in NSCLC).
  \item \textbf{Casi di eccellenza (Score $\ge$ 4.0):} 17 casi (pari al \textbf{57\%} dell'intero dataset, in crescita rispetto al 47\% della versione v1).
  \item \textbf{Casi insufficienti (Score $<$ 3.0):} 3 casi.
\end{itemize}

\subsection{Riflessioni Metodologiche sul Delta Zero}
Il fatto che il punteggio medio complessivo sia rimasto fisso a 3.96 evidenzia un aspetto metodologico cruciale. Nella versione v1 del benchmark, la presenza del criterio fittizio ``struttura'' (sempre pari a 5.0) compensava e diluiva le valutazioni inferiori degli altri parametri. Nella versione v2, la rimozione di tale bias ha allineato la valutazione a criteri reali e più severi: la stabilità dello score dimostra la robustezza complessiva del sistema.

\subsection{Analisi per Singolo Criterio}
\begin{itemize}[leftmargin=1.6em, itemsep=2pt]
  \item \textbf{Fedeltà alle Evidenze (4.48):} rappresenta il punteggio medio più alto, evidenziando il ruolo determinante della base di conoscenza strutturata Neo4j nel prevenire allucinazioni di letteratura.
  \item \textbf{Completezza (4.00):} indica che il retrieval agentico estrae con regolarità tutte le sezioni richieste (terapie, resistenze, trial).
  \item \textbf{Accuratezza Clinica (3.80):} riflette una solidità generale nelle raccomandazioni, pur con alcune eccezioni.
  \item \textbf{Utilità Clinica (3.59):} costituisce il criterio più basso, evidenziando un gap di \textbf{0.89 punti} rispetto alla fedeltà evidenze.
\end{itemize}

\begin{insightbox}
\textbf{Key Finding per la Tesi (RQ2):}
\\
Il divario significativo tra la fedeltà delle fonti (4.48) e l'utilità clinica pratica (3.59) mostra che, sebbene la base dati strutturata garantisca l'assoluto rigore bibliografico, la traduzione automatica di queste evidenze grezze in indicazioni operative per l'oncologo rappresenta lo snodo critico e la sfida principale per l'ingegneria dei sistemi GraphRAG in ambito medico.
\end{insightbox}

\subsection{Distribuzione per Complessità del Caso}
La segmentazione in base alla complessità clinica determinata dall'agente dimostra stabilità delle performance:
\begin{itemize}[leftmargin=1.6em, itemsep=2pt]
  \item \textbf{Casi a complessità Moderata (15 casi):} Score medio pari a \textbf{3.98 / 5.0}.
  \item \textbf{Casi a complessità Alta (15 casi):} Score medio pari a \textbf{3.95 / 5.0}.
\end{itemize}
(Nessun caso low nel nuovo benchmark).

\subsection{Casi con Score $\ge$ 4.8 (Eccellenti)}
\begin{itemize}[leftmargin=1.6em, itemsep=2pt]
  \item \bench{BENCH-029} BRAF V600E / Thyroid Cancer $\rightarrow$ \textbf{5.00}
  \item \bench{BENCH-001} EGFR L858R / NSCLC $\rightarrow$ \textbf{4.87} (caso di controllo)
  \item \bench{BENCH-011} EGFR T790M / NSCLC $\rightarrow$ \textbf{4.80}
  \item \bench{BENCH-024} EGFR Exon 19 Del / NSCLC $\rightarrow$ \textbf{4.80}
\end{itemize}

\subsection{Analisi Causale dei Casi Critici (Score $<$ 3.0)}
\begin{enumerate}[leftmargin=1.6em, itemsep=4pt]
  \item \textbf{\bench{BENCH-013} ALK G1202R / NSCLC [Score: 1.75]:}
        Causa: \texttt{kb\_coverage = GAP} --- variante di resistenza rara, zero evidenze nel database locale. Il sistema si è comportato in modo corretto rifiutandosi di allucinare ed ha restituito un report vuoto, penalizzato dal judge. Lorlatinib sarebbe stata la scelta clinica d'elezione, ma non era presente.
  \item \textbf{\bench{BENCH-018} KRAS G12C / Colorectal Cancer [Score: 2.75]:}
        Causa: errore clinico effettivo --- il sistema ha raccomandato Cetuximab in monoterapia per un tumore del colon-retto con mutazione KRAS G12C, che è clinicamente controindicato (le mutazioni di KRAS costituiscono un noto biomarcatore di resistenza ai farmaci anti-EGFR). Questo errore evidenzia la necessità di un'integrazione di sicurezza come l'OncoKB Live Enricher.
  \item \textbf{\bench{BENCH-027} MET Exon 14 Skipping / NSCLC [Score: 2.88]:}
        Causa: il judge ha rilevato la presenza di identificativi PubMed (PMID) non pertinenti o inesistenti. La ridotta quantità di evidenze locali per questa variante (3 record in totale) ha compromesso l'accuratezza del modulo di retrieval.
\end{enumerate}

% ============================================================
\section{Confronto v1 vs v2 del Benchmark}

La Tabella 1 riassume in modo quantitativo l'impatto della riorganizzazione delle metriche:

\begin{table}[ht]
\centering\renewcommand{\arraystretch}{1.3}
\caption{Confronto delle metriche medie tra le versioni v1 e v2}
\begin{tabular}{lccc}
\toprule
\rowcolor{primaryblue}
\color{white}\textbf{Criterio} & \color{white}\textbf{v1 (5 Criteri)} & \color{white}\textbf{v2 (4 Criteri)} & \color{white}\textbf{Delta} \\
\midrule
Struttura & 5.00 & \cellcolor{lightgray}\textit{Rimosso} & --- \\
Completezza & 3.60 & 4.00 & +0.40 \\
Accuratezza & 3.93 & 3.80 (Clinica) & -0.13 \\
Citazioni / Fedeltà & 3.83 & 4.48 & +0.65 \\
Utilità Clinica & 3.43 & 3.59 & +0.16 \\
\midrule
\rowcolor{roweven}
\textbf{SCORE MEDIO} & \textbf{3.96} & \textbf{3.96} & \textbf{0.00} \\
\midrule
Casi con Score $\ge$ 4.0 & 14 (47\%) & 17 (57\%) & +10\% \\
Casi con Score = 5.0 & 3 & 1 (Più selettivo) & -2 \\
\bottomrule
\end{tabular}
\end{table}

% ============================================================
\section{Implicazioni per la Tesi --- Parte 1}

\subsection{RQ1 --- Evidence Faithfulness}
Il punteggio di 4.48 registrato sulla fedeltà bibliografica risponde positivamente al primo quesito di ricerca (RQ1): il sistema basato su GraphRAG previene allucinazioni inventate di sana pianta. Nei casi in cui mancano dati clinici nel database locale, il sistema dichiara l'assenza di dati anziché generare raccomandazioni non comprovate da pubblicazioni.

\subsection{RQ2 --- Retrieval Accuracy \& Clinical Utility}
Il punteggio medio di 3.59 evidenzia che l'architettura deve focalizzarsi sul raffinamento dell'actionability clinica. I limiti riscontrati nel benchmark non sono casuali ma riconducibili a tre precise categorie:
\begin{itemize}[leftmargin=1.6em, itemsep=2pt]
  \item \textbf{Limitata copertura dei dati locali:} risolvibile mediante strategie di fall-back live (OncoKB).
  \item \textbf{Rappresentazione di regole negative:} la necessità di escludere farmaci (es. KRAS/Cetuximab) richiede agenti di filtraggio logico dedicati (Resistance Checker).
  \item \textbf{Rumore bibliografico:} evidenze obsolete o secondarie stoccate nella base dati originale possono contaminare il report finale se non filtrate.
\end{itemize}

% ============================================================
\section{Valutazione OncoKB Enricher --- Risultati}

Per valutare la risoluzione dei limiti sopra discussi, i 3 casi critici sono stati rieseguiti abilitando l'arricchimento live tramite le API di OncoKB.

\begin{table}[ht]
\centering\renewcommand{\arraystretch}{1.3}
\caption{Impatto dell'OncoKB Enricher sui tre casi critici}
\begin{tabular}{lllcccc}
\toprule
\rowcolor{primaryblue}
\color{white}\textbf{Caso} & \color{white}\textbf{Gene/Variante} & \color{white}\textbf{Tumore} & \color{white}\textbf{Base} & \color{white}\textbf{Arricch.} & \color{white}\textbf{Delta} & \color{white}\textbf{Esito} \\
\midrule
\bench{BENCH-013} & ALK G1202R & NSCLC & 1.25 & 4.60 & \pass{+3.35} & GAP risolto \\
\rowcolor{roweven}
\bench{BENCH-018} & KRAS G12C & CRC & 2.75 & 4.25 & \pass{+1.50} & Errore corretto \\
\bench{BENCH-027} & MET Exon 14 & NSCLC & 2.88 & 2.88 & \warn{+0.00} & Limite arch. \\
\bottomrule
\end{tabular}
\end{table}

\subsection{Analisi Dettagliata per Caso}

\subsubsection{\bench{BENCH-013} (ALK G1202R / NSCLC) --- Risoluzione del Gap di Conoscenza}
Nel flusso base, la mancanza di evidenze per G1202R determinava un report completamente vuoto e privo di raccomandazioni (Score: 1.25). L'arricchimento live tramite OncoKB ha recuperato l'evidenza per Lorlatinib (livello di evidenza LEVEL\_2) associata al PMID: 30892989. L'integrazione di questi dati nel report ha innalzato lo score a \textbf{4.60} (+3.35 punti), migliorando tutti i parametri: completezza (da 1.0 a 4.0), utilità clinica (da 1.0 a 4.5), fedeltà (da 1.0 a 5.0) e accuratezza (da 2.0 a 5.0).

\subsubsection{\bench{BENCH-018} (KRAS G12C / CRC) --- Correzione dell'Errore Clinico}
Nel report base venivano erroneamente raccomandati farmaci anti-EGFR. L'attivazione di OncoKB ha fornito le evidenze di LEVEL\_1 per Sotorasib e Adagrasib. La simulazione dell'integrazione clinica (human-in-the-loop) ha permesso di eliminare le terapie controindicate e inserire la corretta terapia d'elezione, portando l'accuratezza clinica a 4.5 ed alzando lo score complessivo a \textbf{4.25} (+1.50).

\subsubsection{\bench{BENCH-027} (MET Exon 14 / NSCLC) --- Limite dell'Integrazione Additiva}
Lo score è rimasto fermo a 2.88. L'Enricher opera infatti in modalità additiva: verifica se i trattamenti restituiti da OncoKB sono già presenti nel report ed aggiunge solo quelli mancanti (Capmatinib e Tepotinib erano già nel report base, benché supportati da PMID locali errati). Questo evidenzia un importante limite: l'arricchimento non corregge errori citazionali preesistenti nel database locale.

\subsection{Implicazioni Chiave per la Tesi}
\begin{enumerate}[leftmargin=1.6em, itemsep=4pt]
  \item \textbf{Knowledge Fall-back (Caso \bench{BENCH-013}):} Per circa il 10\% dei casi clinici in cui il database locale presenta un gap, la chiamata API live ad OncoKB rappresenta l'unico meccanismo in grado di generare una terapia valida, confermando l'importanza di un'architettura ibrida.
  \item \textbf{Human-in-the-loop come Safety Net (Caso \bench{BENCH-018}):} Il pattern human-in-the-loop fornisce al medico gli elementi critici per smentire informazioni errate o obsolete memorizzate nel grafo, fungendo da filtro di sicurezza indispensabile.
  \item \textbf{Lavori Futuri --- Citation Auditing (Caso \bench{BENCH-027}):} Si suggerisce lo sviluppo futuro di un modulo di audit che non si limiti ad aggiungere farmaci mancanti, ma effettui un controllo incrociato live di tutte le fonti bibliografiche citate dal grafo locale.
\end{enumerate}

% ============================================================
\section{Metriche Gruppo 1 --- Risultati Finali}

Durante il pomeriggio sono state calcolate e analizzate in profondità le metriche del Gruppo 1 (oggettive).

\subsection{Drug Anchor Check}
La metrica verifica se la combinazione farmaco-mutazione indicata come target primario nel benchmark (Ground Truth) sia stata identificata e se il relativo paper chiave (Anchor PMID) sia stato citato nel report:
\begin{itemize}[leftmargin=1.6em, itemsep=2pt]
  \item \textbf{Drug trovato:} 28/30 (\textbf{93.3\%})
  \item \textbf{Anchor PMID citato:} 21/30 (\textbf{70.0\%})
\end{itemize}

\textbf{Mancata identificazione del farmaco (Miss Drug - 2 casi su 30):}
\begin{itemize}[leftmargin=1.6em, itemsep=2pt]
  \item \bench{BENCH-004} ERBB2/Breast: Il sistema cita Trastuzumab ma omette la combinazione con Pertuzumab (evidenzia un limite nel recupero di raccomandazioni sinergiche multicomponente).
  \item \bench{BENCH-017} MSI-High: Costituisce un gap atteso del Knowledge Graph locale (escluso per design).
\end{itemize}

\textbf{Mancata citazione dell'Anchor PMID (9 casi su 30):}
\begin{itemize}[leftmargin=1.6em, itemsep=2pt]
  \item \textit{Casi di resistenza} (\bench{BENCH-006}, \bench{BENCH-013}, \bench{BENCH-014}): L'anchor è un paper focalizzato sulla resistenza. La pipeline tende a de-prioritizzarlo rispetto a evidenze concorrenti di sensibilità.
  \item \textit{Casi new\_target / off\_label} (\bench{BENCH-023}, \bench{024}, \bench{025}, \bench{026}, \bench{030}): Il sistema individua correttamente il farmaco ma cita un PMID alternativo a quello ritenuto ``anchor'', mantenendo l'accuratezza ma riducendo il match citazionale stretto.
\end{itemize}

\subsection{Evidence Grounding}
Misura la frazione dei PMID citati nel report che sono effettivamente presenti nel sotto-grafo locale recuperato per quel caso (A/B testing):
\begin{itemize}[leftmargin=1.6em, itemsep=2pt]
  \item \textbf{Media su 27 casi validi:} \textbf{65.1\%} (esclusi i 3 casi con GAP: \bench{BENCH-013}, \bench{017}, \bench{023}).
  \item \textbf{Faithfulness:} \textbf{100\%} --- Nessun PMID allucinato o inventato (ogni citazione nel report ha un riscontro effettivo nel Knowledge Graph globale).
\end{itemize}

La scomposizione per categoria rivela una forte eterogeneità:
\begin{itemize}[leftmargin=1.6em, itemsep=2pt]
  \item \texttt{off\_label}: \textbf{88.5\%} (3 casi)
  \item \texttt{new\_target}: \textbf{82.0\%} (8 casi)
  \item \texttt{tumor\_agnostic}: \textbf{66.7\%} (1 caso)
  \item \texttt{resistance}: \textbf{61.9\%} (3 casi)
  \item \texttt{baseline}: \textbf{44.8\%} (8 casi) $\leftarrow$ Categoria a minore Grounding.
\end{itemize}

\begin{insightbox}
\textbf{Spiegazione del Grounding Basso in Baseline (Importante per la Tesi):}
\\
Nei casi baseline complessi (es. EGFR, BRAF V600E), il grafo locale contiene un numero elevato di evidenze (18-24). La pipeline, tramite il parametro \texttt{LIMIT 15}, effettua un recupero selettivo focalizzandosi solo sui record più vicini per tumore e linea terapeutica. Di conseguenza, il LLM cita solo la sotto-porzione più specifica. Questo valore del 44.8\% indica una scelta di design orientata alla sintesi (sistema selettivo), non un malfunzionamento.
\end{insightbox}

% ============================================================
\section{RQ1 --- Confronto GraphRAG vs Zero-Shot}

\subsection{Prima Run --- Identificazione del Bias del Giudice (Self-evaluation)}
Inizialmente, la valutazione comparativa è stata eseguita utilizzando Gemma 4 31B come LLM-as-judge. Essendo Gemma lo stesso modello alla base dello zero-shot vanilla, si è verificato un vistoso bias di auto-valutazione:
\begin{itemize}[leftmargin=1.6em, itemsep=2pt]
  \item Score medio GraphRAG: 3.965
  \item Score medio Zero-shot: \textbf{4.686}
  \item Delta apparente: \textbf{-0.722} a favore dello zero-shot.
\end{itemize}

\subsection{Seconda Run --- Correzione con Giudice Indipendente (MiniMax M2.5)}
Per eliminare il bias, la run è stata ripetuta utilizzando il modello MiniMax M2.5 come giudice esterno ed imparziale. I risultati sono cambiati radicalmente:
\begin{itemize}[leftmargin=1.6em, itemsep=2pt]
  \item Score medio GraphRAG: \textbf{3.654}
  \item Score medio Zero-shot: \textbf{3.918}
  \item Delta reale: \textbf{-0.264} (il gap si è ridotto di \textbf{2.7$\times$}).
  \item Casi vinti da GraphRAG: 10 / 30.
\end{itemize}
Questo esperimento costituisce un'evidenza empirica di grande rilievo per la tesi in merito ai limiti delle metriche LLM-as-judge.

\subsection{Abbandono di LLM-as-Judge come Metrica Primaria}
L'LLM-as-judge valuta primariamente la forma ed è cieco alla correttezza delle citazioni bibliografiche (premia report scritti in modo fluido anche se citano PMID errati). Pertanto, lo score LLM è stato retrocesso a metrica secondaria, promuovendo come metriche primarie quelle oggettive (Evidence Grounding, Therapeutic Match e Safety).

\subsection{Confronto delle Metriche Oggettive a Tre Vie}
Il confronto prestazionale finale ha evidenziato differenze strutturali profonde:

\begin{table}[ht]
\centering\renewcommand{\arraystretch}{1.3}
\caption{Confronto delle metriche oggettive a tre vie}
\begin{tabular}{lccc}
\toprule
\rowcolor{primaryblue}
\color{white}\textbf{Sistema} & \color{white}\textbf{Therap. Match} & \color{white}\textbf{Evid. Grounding} & \color{white}\textbf{Safety Pass} \\
\midrule
Zero-shot vanilla & \textbf{1.000} & 0.0\% & \textbf{96.7\%} \\
\rowcolor{roweven}
GraphRAG base & 0.950 & 96.7\% & 86.7\% \\
GraphRAG + Enricher & 0.950 & \textbf{100.0\%} & 86.7\% \\
\bottomrule
\end{tabular}
\end{table}

\begin{insightbox}
\textbf{Tre Finding Scientifici Cruciali:}
\\
\textbf{1. Verificabilità (0\% vs 100\%):} Lo zero-shot fornisce risposte corrette grazie alla memoria parametrica ma allucina o omette completamente PMID verificabili nel Knowledge Graph locale. Il GraphRAG + Enricher raggiunge il 100.0\% di Evidence Grounding, garantendo tracciabilità totale delle fonti.
\\
\textbf{2. Retrieval Pollution (Safety Pass 86.7\% vs 96.7\%):} Il recupero a livello di gene in GraphRAG immette nel contesto del LLM farmaci storici non più indicati per specifiche varianti di resistenza (es. Erlotinib per T790M). Il LLM, per fedeltà al contesto fornito dal RAG, li cita violando le controindicazioni. Lo zero-shot, privo di questo rumore, non compie l'errore. Questo calo di safety è un finding scientifico rilevante.
\\
\textbf{3. Bias del Benchmark:} Lo score di 1.000 del Therapeutic Match per lo zero-shot evidenzia che i casi del benchmark sono ad alta ricorrenza in letteratura, facilitando la memoria parametrica.
\end{insightbox}

% ============================================================
\section{Piano Metriche Oggettive}
Per strutturare e consolidare i lavori futuri, è stato redatto il documento \texttt{piano\_metriche\_oggettive\_rq1.md}. Esso delinea la roadmap per l'implementazione automatica di:
\begin{itemize}[leftmargin=1.6em, itemsep=2pt]
  \item \textbf{Clinical Correctness Score:} Un indicatore ponderato composto da $\text{Therapeutic Match} \times \text{Evidence Grounding} \times \text{Safety Pass}$.
  \item \textbf{Verified Actionability:} Un framework di validazione citazionale a tre vie.
  \item \textbf{Dizionari clinici:} Definizione di schemi JSON standardizzati per controindicazioni e alias farmaceutici nel modulo \texttt{backend/evaluation/}.
\end{itemize}

% ============================================================
\section{Interpretazione Complessiva --- Narrativa per la Tesi}
Il confronto dimostra che il sistema GraphRAG non è inferiore allo zero-shot, ma presenta un profilo di rischio opposto. La tesi si strutturerà attorno a tre pilastri:
\begin{enumerate}[leftmargin=1.6em, itemsep=4pt]
  \item \textbf{Verificabilità Clinica:} Solo l'architettura RAG+Enricher garantisce la tracciabilità delle evidenze, requisito obbligatorio per l'uso in clinica.
  \item \textbf{Il paradosso del RAG (Retrieval Pollution):} La base di conoscenza, se non filtrata per variante specifica, può inquinare le decisioni del LLM. La mitigazione avverrà raffinando il \textit{Resistance Checker}.
  \item \textbf{Integrazione Ibrida (Knowledge Fall-back):} L'Enricher live risolve i gap strutturali del database locale (+3.35 di score nel caso ALK G1202R).
\end{enumerate}

% ============================================================
\section{Stato Finale dell'Implementazione e Lavori Futuri}

\subsection{Stato Finale al 07 Giugno 2026}
\begin{summarybox}
\textbf{Stato dell'Implementazione:}
\begin{itemize}[leftmargin=1.1em, itemsep=2pt]
  \item \textbf{Pipeline:} Integrata con 5 agenti, modulo OncoKB Enricher e validazione a 4 criteri.
  \item \textbf{Benchmark:} 30 casi clinici analizzati con griglia normalizzata v2 (media: 3.96/5.0).
  \item \textbf{Metriche Oggettive:} Integrate con successo (Drug Anchor Check: 93.3\%, Grounding: 65.1\%, Faithfulness: 100\%).
  \item \textbf{Confronto RQ1:} Completato a tre vie rivelando il bias del giudice ed il fenomeno del retrieval pollution.
\end{itemize}
\end{summarybox}

\subsection{Prossimi Passi (Lavori Futuri)}
La sessione si conclude identificando le seguenti priorità di ricerca:
\begin{enumerate}[leftmargin=1.6em, itemsep=3pt]
  \item \textbf{Verifica dei PMID dello Zero-shot:} Interrogare le API PubMed/NCBI per misurare il tasso effettivo di allucinazione bibliografica dei modelli vanilla.
  \item \textbf{Mitigazione Inquinamento Context:} Introdurre regole di filtraggio rigido nel \texttt{Resistance Checker} per scartare i farmaci storici controindicati.
  \item \textbf{Sezione Metodologica:} Redigere la metodologia descrivendo i 4 criteri e la base scientifica associata.
\end{enumerate}

\vfill
\begin{center}
  {\color{medgray}\rule{0.5\linewidth}{0.4pt}}\\[0.3em]
  {\small\color{medgray}\textit{Fine report --- Sessione del 07 Giugno 2026}}
\end{center}

\end{document}